\doublespacing
\section{Hypotheses}
\subsection{Culture-talk is more common among high-status groups}
Extant theoretical accounts of the role of culture-talk in the status stratification process provide some partial guidance as to what we should expect to find.  First, a direct empirical implication of theories of cultural stratification that isolate the role of ``entertainment talk'' in the interactional lives of members of the cosmopolitan middle class leads us to expect that \emph{the propensity to engage in culture-talk increases with such status markers as education, especially in domains usually marked as ``legitimate culture'' even after adjusting for other confounding factors such as gender, income, age or family status}.\footnote{By ``legitimate culture'' I refer to such domains as the fine and performance arts delivered through non-profit and mixed systems \citep{dimaggio00} or the book-based print culture that has become increasingly concentrated within the mobile ``reading class'' \citep{griswold08}.} 

\subsection{Status linked differences in culture-talk can be found in all cultural domains}
Integrating this line of thinking with recent developments in cultural stratification centered around the emergence of cultural omnivorousness as the now established form of cultural capital among the mobile upper-middle class \citep{ollivier04, katzgerro04, ollivier08, warde_etal08, peterson05, dimaggio96, vaneijck00} leads to a more general implication: \emph{the propensity to engage in culture-talk is higher among the most culturally advantaged status groups across all consumption domains, including those usually classified as ``popular'' (e.g. commercial arts delivered through large scale culture industry systems)}.  Thus, just like the new upper-middle class display omnivorous \emph{tastes}, they should also display an equal-opportunity propensity towards the inclusion of all types of cultural content in interaction (``culture-talk omnivorousness'').

\subsection{Consumption partially mediates status group differences in culture talk across domains}
Cultural consumption should be positively associated with higher rates of culture talk across all domains, since direct familiarity with the relevant symbolic goods is a (partial) precondition for the mobilization of culture in interaction \citep{dimaggio87}. Because high status individuals---e.g. indexed by levels of schooling \citep{bourdieu84, bourdieu_darbel91, meyer77}---already have a higher predisposition to engage in aesthetic consumption \emph{across domains} (as has been repeatedly demonstrated in previous research on the omnivore phenomenon) we should find that \emph{once we adjust for rates of consumption differences across educational groups should be smaller}. This would suggest that consumption is one of the intervening \emph{mechanisms} explaining heterogeneity in culture-talk propensities across status groups. 

\subsection{The mediating effect of consumption is weaker for popular culture}
The mediating effect of consumption in relation to status group differences (if it exists) should not apply with equal strength to all forms of culture-talk. In other words, in cultural domains associated with the popular and commercial arts differential probabilities (linked to status group membership) of engaging in culture-talk may be partially independent from consumption.  The reason for this is simple: participation in popular or commercial culture should not differ much across such status markers as levels of education, while participation in more legitimate culture should be more tightly linked to status group membership.  If member of high-status groups are more likely to engage in culture-talk \emph{across domains} (as intimated above), then adjusting for rates of engagement in cultural domains where participation is more loosely connected to status group membership should not completely eliminate cross-group differences in the likelihood of engaging in culture-talk.  This leads us to expect that for \emph{domains more likely to be characterized as popular culture, status-group differences in the rate of culture talk is partially independent from rates of participation}. Note that linking a test implication derived from micro-interactionist cultural stratification theory and recent research on the omnivore phenomenon (as applied to culture-talk) leads to a somewhat counter-intuitive implication: \emph{high-status individuals should be able to mobilize culture as talk even in consumption domains where they do not enjoy a consumption advantage in relation to low status groups}.\footnote{Note also that if status-linked heterogeneity in culture-talk propensity do not exist in the popular domain, then this will not be the case.} 

\subsection{High-status consumers of legitimate culture are more likely to engage in culture talk in commercial domains}
If it is true that status-linked differences in culture-talk should be partially independent from consumption in domains in which consumption is not correlated with status, then the question becomes what explains culture-talk heterogeneity across status groups in these cultural domains?  In what follows I explore one possible answer to this question: \emph{high-status persons who are culturally active in the domain of legitimate culture are more likely to deploy cultural content in interaction}. A test implication of this proposition is that \emph{once we adjust for the propensity to consume other forms of more legitimate culture, we should find that status-linked differences in culture talk propensities should be largely accounted for}.  This entails that it is culturally active members of the new middle classes who are responsible for status-linked differences in the mobilization of culture in interaction in domains where consumption is independent from status.  

\subsection{Consumption of legitimate culture increases culture talk in popular domains but not vice versa}
The prediction that culturally active members of high-status groups are primarily responsible for the culture-talk advantage of this group in popular culture domains, naturally leads to a substantively important empirical implication: there should exist an asymmetry in the effects of culture consumption on culture talk when we predict culture-talk in ``popular'' domains using observed consumption propensities in more legitimate domains and vice versa, such that \emph{the propensity to engage in high-status culture consumption helps us predict culture talk in popular domains but not the reverse}. That is, knowing that somebody is a fine arts consumer or a member of the reading class should help us predict that she is more likely to engage in culture-talk in commercial and popular domains (holding constant the propensity to consume culture in that domain), but knowing that somebody consumes popular culture should not give us any additional leverage in predicting whether somebody mobilizes fine arts or literary culture in interaction (after adjusting for the propensity to be an active fine arts consumer or an active reader). 
